%% LyX 2.3.6.2 created this file.  For more info, see http://www.lyx.org/.
%% Do not edit unless you really know what you are doing.
\documentclass[english]{article}
\usepackage[T1]{fontenc}
\usepackage[latin9]{inputenc}
\usepackage{bm}
\usepackage{amsmath}
\usepackage{babel}
\begin{document}
Here we start from the effective model described in {[}SUNKO{]} and
derive the effective EME coupling that is presented in the main text.
The key new insight is as follows: the effective Kondo-like coupling
that is generated between the Pd-electrons and Cr local moments involves
an inter-layer hybridization, $g_{ij}$ between sites $i$ and $j$.
If this coupling were to be affected by the fluctuations associated
with the corresponding out-of-plane bond, $g_{ij}$ will be renormalized
by the displacement, $\varphi_{ij}$. Following similar steps as in
{[}SUNKO{]} we define the creation and annihilation operators, $p_{i}^{\left(\dagger\right)},c_{i}^{\left(\dagger\right)}$and
for the Pd an Cr electrons respectively. With these defintions, the
Hamiltonian of the system is given by:

\begin{align}
H & =H_{\text{Pd}}+H_{\text{Cr}}+H_{\text{hyb}}+H_{\text{ph,EME}},\\
H_{\text{Pd}} & =\sum_{\boldsymbol{k}\sigma}\left(\varepsilon_{\boldsymbol{k}}-\mu\right)p_{\boldsymbol{k}\sigma}^{\dagger}p_{\boldsymbol{k}\sigma},\\
H_{\text{Cr}} & =-\sum_{ij\sigma}t_{ij}^{c}c_{i}^{\dagger}c_{j}+\text{h.c}+U\sum_{i}\left(n_{i\uparrow}^{c}-\frac{1}{2}\right)\left(n_{i\downarrow}^{c}-\frac{1}{2}\right),\\
H_{\text{hyb}} & =\sum_{\left\langle ij\right\rangle \sigma}g_{ij}\left(1-\lambda\varphi_{ij}\right)\left(c_{i\sigma}^{\dagger}p_{j\sigma}+\text{h.c}\right)\\
H_{\text{ph,EME}} & =\sum_{\left\langle ij\right\rangle }\frac{\pi_{ij}^{2}}{2M}+\frac{M\tilde{\omega}_{0}}{2}\varphi_{ij}^{2}.
\end{align}

here $t_{ij}^{c}$ is the hopping in the Cr layer, $\varepsilon_{\boldsymbol{k}}$
is the dispersion for th Pd electrons, where we keep hopping up to
second nearest neighbours, $\lambda$ determines the scale by which
the hybridization is renormalized to leading order by the bond displacement
$\varphi_{ij}$, $M$ is the phonon mass, $n_{i\sigma}^{c}$ is the
density operator for Cr electrons, and $U$ is the on site interaction
strenght. In general these interactions would involve multiple orbitals
in the Cr layer and contain the full lattice structure of PdCrO$_{2}$,
for simplicitly, however, we only consider a single orbital per site
in the Cr layers.

Since it has been demonstrated that the Cr atoms are strongly correlated
and form a Mott insulator with 120$^{\circ}$ AFM order at low temperatrures,
it is justified to take the strong coupling limit $U\gg t^{c},g$
to get a low energy effective theory in terms of local moments in
teh Cr layer {[}SUNKO{]}. We can obtain this low energy theory by
performing a Schriffer-Wolf transformation of the above Hamiltonian.
If we only keep linear terms in the bond displacement, the resulting
Hamiltonian will be given by :

\begin{align}
H & =H_{\text{Pd}}+H_{\text{lm}}+H_{\text{K}}+H_{\text{EME}}+H_{\text{ph,EME}},\\
H_{\text{lm}} & =J_{H}\sum_{\left\langle ij\right\rangle }\boldsymbol{S}_{i}\cdot\boldsymbol{S}_{j}\\
H_{\text{K}} & =\sum_{\left\langle ilj\right\rangle \sigma}\frac{4g_{ij}g_{lj}}{U}\boldsymbol{S}_{l}\cdot\left(p_{i}^{\dagger}\boldsymbol{\sigma}p_{j}\right)\\
H_{\text{EME}} & =-\sum_{\left\langle ilj\right\rangle \sigma}\frac{4g_{ij}g_{lj}}{U}\lambda\left(\varphi_{il}+\varphi_{jl}\right)\boldsymbol{S}_{l}\cdot\left(p_{i}^{\dagger}\boldsymbol{\sigma}p_{j}\right)
\end{align}

The first term above is the usual Heisenberg coupling with antiferromagnetic
exchange $J_{H}=4t^{c}/U,$ where we only considered nearest neighbour
hopping between Cr sites. The second term above is the Kondo coupling
that was reported in {[}SUNKO{]}, while the third is the new EME coupling
between local moments in the Cr layer,optical phonons and itinerant
electrons. In what follows we consider the simplified case where adjacent
Pd layers are aligned completely aligned and the hibridization is
a constant independent of bond direction $g_{ij}\rightarrow g$. In
this limit, we can define the Kondo and EME couplings to be $J_{K}=4g^{2}/U$
and $\tilde{\alpha}=-J_{K}\lambda$. We then have the following Hamiltonian
in momentum space
\begin{equation}
H=H_{\text{Pd}}+H_{\text{ph}}+H_{\text{lm}}+H_{\text{int}}
\end{equation}

with the interaction given by
\begin{align}
H_{\text{int}}= & \frac{1}{\sqrt{N}}\sum_{\substack{\bm{k}\bm{k}'\\
\alpha\beta
}
}J_{K}\left(\boldsymbol{k},\boldsymbol{k}'\right)p_{\boldsymbol{k}\alpha}^{\dagger}\left[\boldsymbol{S}_{\boldsymbol{k}-\boldsymbol{k}'}\cdot\boldsymbol{\sigma}_{\alpha\beta}\right]p_{\boldsymbol{k}'\beta}\\
 & +\frac{1}{N}\sum_{\bm{k}\bm{k}'\bm{q}}\sum_{\ell\alpha\beta}\tilde{\alpha}_{\ell}\left(\boldsymbol{k},\boldsymbol{k}'\right)\varphi_{\ell}\left(\boldsymbol{k}-\boldsymbol{k}'-\boldsymbol{q}\right)p_{\boldsymbol{k}\alpha}^{\dagger}\left[\boldsymbol{S}_{\boldsymbol{q}}\cdot\boldsymbol{\sigma}_{\alpha\beta}\right]p_{\boldsymbol{k}'\beta}
\end{align}

where the Fourier transfomed operators are:
\begin{equation}
\varphi_{\ell+j,j}=\frac{1}{\sqrt{N}}\sum_{\bm{k}}e^{-i\bm{k}\bm{R}_{j}}\varphi_{\ell}\left(\boldsymbol{k}\right)
\end{equation}

\begin{equation}
p_{i\alpha}^{\dagger}=\frac{1}{\sqrt{N}}\sum_{\bm{k}}e^{-i\bm{k}\bm{R}_{i}}p_{\boldsymbol{k}\alpha}^{\dagger}
\end{equation}
\begin{equation}
\boldsymbol{S}_{j}=\frac{1}{\sqrt{N}}\sum_{\bm{q}}e^{-i\bm{q}\bm{R}_{j}}\boldsymbol{S}_{\boldsymbol{q}}
\end{equation}

and the momentum dependent coupligns are 
\begin{align*}
J_{K}\left(\boldsymbol{k},\boldsymbol{k}'\right) & =J_{K}\left[2\cos\left(k_{z}a_{z}\right)\right]\left[2\cos\left(k_{z}'a_{z}\right)\right]
\end{align*}

and

\begin{align*}
\tilde{\alpha}_{+1}\left(\boldsymbol{k},\boldsymbol{k}'\right) & =\tilde{\alpha}_{-1}\left(\boldsymbol{k},\boldsymbol{k}'\right)=\tilde{\alpha}\left(\boldsymbol{k},\boldsymbol{k}'\right)=-\lambda J_{K}\left(\boldsymbol{k},\boldsymbol{k}'\right)
\end{align*}

with $\bm{R}_{i}$ being the position of site $i$ and $N$ the number
of unit cells as the normalization of the Fourier transform. 

In the above derivation we introduced several three new parameters
to the theory in {[}SUNKO{]}, the mass of the EME mode $M$, its frequency
$\tilde{\omega}_{0}$ and the EME coupling $\lambda$. As we detail
in the text, these parameters only enter the theory either in the
precise combination that corresponds to the dimensionless EME coupling,
or in the combination $\beta\tilde{\omega}_{0}$, leaving effectively
two new parameters in the theory. The rest of the parameters are taken
from the values provided in {[}SUNKO{]}. Specifically, we have $U=4$eV
and $g=0.1$eV which lead to a Kondo coupling of $J_{K}=12$ meV.
Additionally, we have $J_{H}=6.25$meV, $t_{nn}^{p}=568$meV, $t_{nnn}^{p}=108$meV,
and $\mu=256$meV. Assuming these values we then fit the difference
of the resistivity of PdCrO$_{2}$ and PdCoO$_{2}$ to obtain the
dimensionless EME coupling in the main text and the frequency of this
mode. 
\end{document}
